% Generated draft: requires mathematical and source review.
\documentclass[11pt]{article}
\usepackage[margin=25mm]{geometry}
\usepackage{fontspec}
\setmainfont{Latin Modern Roman}
\usepackage{amsmath,amssymb,mathtools,unicode-math}
\setmathfont{Latin Modern Math}
\usepackage{enumitem,needspace,xurl,hyperref}
\setlength{\parindent}{0pt}
\setlength{\parskip}{5pt}
\setlength{\emergencystretch}{3em}
\newcommand{\N}{\mathbb{N}}
\newcommand{\Z}{\mathbb{Z}}
\newcommand{\Q}{\mathbb{Q}}
\newcommand{\R}{\mathbb{R}}
\newcommand{\C}{\mathbb{C}}
\newcommand{\F}{\mathbb{F}}
\newcommand{\A}{\mathbb{A}}
\newcommand{\PP}{\mathbb{P}}
\newcommand{\E}{\mathbb{E}}
\newcommand{\eps}{\varepsilon}
\newcommand{\dd}{\,\mathrm{d}}
\begin{document}
\section*{UnsolvedMath ID: 30001895}
\textbf{Source problem code:} catalogue entry
\subsection{Definitions and mathematical statement}
Let \(p,q\) be positive integers with \(q\le p\). Under the usual common-intersection interpretation, a family \(\mathcal F\) of sets has the \((p,q)\)-property if
\[
\forall\,\mathcal G\subseteq\mathcal F,\quad
|\mathcal G|=p\ \Longrightarrow\
\exists\,\mathcal H\subseteq\mathcal G:\quad
|\mathcal H|=q\quad\text{and}\quad
\bigcap_{A\in\mathcal H}A\ne\varnothing.
\]
For a family of subsets of an ambient set \(X\), a point transversal is a set \(T\subseteq X\) such that \(T\cap A\ne\varnothing\) for every \(A\in\mathcal F\). Interpret a \(k\)-transversal as such a set with \(|T|\le k\). These interpretations of the terminology require confirmation against the source report, since the supplied record does not define them.

The geometric question is the following proposed implication. Given a family \(\mathcal F\) of hyperplanes in \(\R^d\), and integers satisfying \(d+1\le q\le p\), does
\[
\mathcal F\text{ has the }(p,q)\text{-property}
\quad\Longrightarrow\quad
\exists\,T\subseteq\R^d:\quad
|T|\le p-q+1\quad\text{and}\quad
\forall H\in\mathcal F,\ T\cap H\ne\varnothing
\]
always hold? Here a hyperplane is provisionally interpreted as an affine hyperplane
\[
H=\{x\in\R^d:\langle a,x\rangle=b\},
\qquad a\in\R^d\setminus\{0\},\quad b\in\R.
\]
The supplied record does not explicitly distinguish affine from linear hyperplanes.

The set-theoretic question asks for the analogous implication. Given an ambient set \(X\), a family \(\mathcal A\subseteq\{A\subseteq X:|A|=r\}\), and integers satisfying \(r+1\le q\le p\), does
\[
\mathcal A\text{ has the }(p,q)\text{-property}
\quad\Longrightarrow\quad
\exists\,T\subseteq X:\quad
|T|\le p-q+1\quad\text{and}\quad
\forall A\in\mathcal A,\ T\cap A\ne\varnothing
\]
always hold?

\emph{Unresolved scope qualifications.} The record does not specify whether either family must be finite, whether it must have at least \(p\) members, or whether families are indexed and may contain repeated members. With the displayed definition, the \((p,q)\)-property is vacuous when a family has fewer than \(p\) members; any intended size hypothesis must therefore be recovered from the report. No such hypothesis has been silently added. The displays state the proposed bounds for a given family; the intended universal scope over families remains subject to these qualifications. The record also does not explicitly state the lower boundary conventions for \(d\) and \(r\), or whether a \((p-q+1)\)-transversal means at most or exactly that many points. This is a statement-only draft, not an assertion that either implication is true.
\subsection{Short English statement}
The supplied problem asks whether a family in which every \(p\) members contain \(q\) members with a common point can always be pierced by at most \(p-q+1\) points: for hyperplanes when \(d+1\le q\le p\), and for \(r\)-element sets when \(r+1\le q\le p\). The intended family-size assumptions and terminology need confirmation. The supplied open label is not independently verified.
\subsection{Sources}
\begin{itemize}
\item Supplied local UnsolvedMath registry record, permanent ID \(30001895\), title \emph{Exact Transversals for Families with the \((p,q)\)-Property}, collection \emph{Oberwolfach Reports Open Problems}, exact source problem code \texttt{catalogue entry}. The record gives proposed year 2011 and names no individual proposer.
\item Supplied report citation: \emph{Discrete Geometry (2012)}, Oberwolfach Reports, pp.\ 2459--2548, DOI \texttt{10.4171/owr/2011/44}. Supplied source URL and literature-source link labeled ``OWR 44 (2012)'': \url{https://doi.org/10.4171/owr/2011/44}. The background describes the reports as documenting mathematical workshops and their open-problem discussions.
\item Supplied literature-source link labeled ``Alon--Kalai variant'': \url{https://arxiv.org/abs/1409.1194}. No fuller bibliographic identification or explanation of its precise applicability is supplied.
\item Literature assessment dated \texttt{2026-08-21}: the supplied curation labels the problem open and states that the exact \(p-q+1\) transversal assertion is not established in the general Hadwiger--Debrunner range, while later work provides finiteness or bounds rather than the universal exact value. This is supplied dated curation derived from the OWR material, not an independently verified literature assessment. The supplied background expressly calls for independent verification of time-sensitive status and cited claims.
\item The registry labels the statement \texttt{corrected\_verified}, but its accompanying verification note states only that TeX notation was enclosed in math delimiters without changing mathematical content. This records a formatting correction, not an independent mathematical verification.
\end{itemize}
\end{document}
